\par\bigskip\SourceLeaf{163}
\Needspace{8\baselineskip}\pdfbookmark[1]{III. Kommissionsbericht}{witwen-III}
\begin{center}
[III.]\par\medskip
[\textit{Kommissionsbericht.}]
\end{center}
\noindent\hspace*{\parindent}Die für die Angelegenheit der Witwenkasse ernannte Kommission beehrt sich, hiermit das Ergebnis ihrer Beratung vorzulegen.

Zuerst stellte sich die Frage, welcher Zinssatz den zu treffenden Maßnahmen zugrunde gelegt werden soll.

Zwar ist das Kapital der Witwenkasse gegenwärtig größtenteils noch zu einem höheren Zinssatz als \(3\tfrac{1}{2}\) Prozent angelegt; bei angemessener Würdigung der Verhältnisse muss der durchschnittliche Zinssatz jedoch auf jeden Fall als unter 4 Prozent liegend betrachtet werden. In Anbetracht der folgenden Gründe:

1) dass der Zinssatz in allen Ländern Europas trotz einzelner Schwankungen ganz unverkennbar die Tendenz zeigt, weiter zu sinken\footnote{Einer der aufgeklärtesten Staatsmänner, namentlich in Finanzfragen, der großherzoglich badische Minister \textsc{Nebenius}, sagt in seinem bekannten Werk Über die Herabsetzung der Zinsen der öffentlichen Schulden, S. 129: ‘Das allmähliche weitere Sinken des Zinssatzes, das bei länger anhaltendem Frieden nicht ausbleiben kann, wird zuletzt überall, hier etwas früher, dort etwas später, die Senkung auf drei Prozent herbeiführen oder sie wenigstens als eine Maßnahme erscheinen lassen, die nur vom Entschluss der Regierung abhängt’; und an anderer Stelle, S. 21: ‘Auf eine Zeit größerer Betriebsamkeit bei produktiven Unternehmungen, die das Sinken des Zinssatzes eine Zeit lang aufhält, folgt umso gewisser ein rasches Sinken des Zinssatzes.’};

2) dass jede wesentlich vom Zinssatz abhängige Anstalt, wenn sie nicht als gänzlich unsicher gelten soll, auf einem etwas niedrigeren als dem augenblicklich bestehenden Zinssatz beruhen muss;

3) dass uns auch die Regelungen von 1794 insofern mit gutem Beispiel vorangegangen sind, als damals wenigstens die Absicht bestand, die Kasse für einen Zinssatz von 3 Prozent abzusichern;

\noindent hält es die Kommission nicht für zulässig, einen höheren Zinssatz als \(3\tfrac{1}{2}\) Prozent zugrunde zu legen. Deshalb hält sie es für unumgänglich notwendig, sich die Deckung des bei diesem Zinssatz entstehenden Defizits von 17520 Talern zum Ziel zu setzen.

Glaubte man, die bei den bisherigen Einrichtungen gegenwärtig noch entstehenden jährlichen Überschüsse {\renewcommand{\thefootnote}{**)}\footnote{Die Jahresrechnung 1844—1845 ergibt keinen Überschuss; denn
\par\noindent für 1844 Juli 1 war das Geldvermögen mit 116375 Thl. 14 Ggr. 6 Pf. ausgewiesen,
\par\noindent für 1845 Juli 1 dagegen nur mit . . . . . . . . . . . . . . . . . 116369 ,, 17 ,, 4 ,,
\par\noindent Da es im laufenden Jahr jedoch drei Witwen weniger gibt, ist in diesem Jahr wieder mit einem Überschuss zu rechnen, der wahrscheinlich auch noch einige Zeit anhalten wird.}} seien schon an sich Schritte zur Deckung des Defizits, und dieses werde getilgt sein, sobald sich das jetzige Vermögen um 17520 Taler vergrößert habe, so würde eine solche Meinung nur auf einer Begriffsverwirrung und einem gänzlichen Verkennen der Bedeutung des Defizits beruhen. Im Ergebnis der Rechnung, mit der die Größe des Defizits ermittelt wurde, sind die gegenwärtig noch entstehenden jährlichen Ersparnisse dem Sinn dieser Rechnung nach bereits vollständig mitberücksichtigt. Diese als Schritte zur Deckung anzusehen, wäre dasselbe, als würde man sie zweimal in Rechnung stellen {\renewcommand{\thefootnote}{***)}\footnote{Da dem Vernehmen nach hier und da ähnliche Ansichten geäußert worden sind, dieser Bericht aber nicht der Ort zu ihrer Widerlegung wäre, ist der Anlage noch eine nähere Beleuchtung des Gegenstandes beigefügt worden. Für alle, die die Denkschrift schon mit der nötigen Aufmerksamkeit gelesen und erwogen haben, wird sie allerdings ganz überflüssig sein.}}. Die Genossenschaft möge sich also keine Illusion über die Wahrheit
\SourceLeaf{164}
machen: Die Deckung des Defizits kann und muss einzig und allein durch zweckmäßige Änderungen der bisherigen Einrichtungen bewirkt werden.

Es versteht sich von selbst, dass schwache Mittel auch nur schwache Wirkungen hervorbringen können. Dazu würde gehören, die bisher freiwillige Teilnahme für alle künftig ernannten Professoren verpflichtend zu machen und die Freiheit aufzuheben, jederzeit wieder auszutreten. Dass eine solche Maßnahme nur eine sehr geringe Wirkung haben könnte, zeigt der Umstand, dass nach einem Durchschnitt über 50 Jahre die mittlere Anzahl der nicht beitragenden, obwohl zur Teilnahme berechtigten Professoren nur \(=4{,}4\) betrug. Die dadurch der Kasse bei bisheriger Beitragshöhe jährlich entgehende Einnahme betrug also \(=44\) Taler. Dies entspricht somit je nach dem Zinssatz einem dauernden Kapital von 1100 oder 1257 Rthl. Dies wird jedoch, wenn auch nicht ganz, so doch größtenteils durch die Strafgelder aufgewogen, die nach der bestehenden, sehr zweckmäßigen Einrichtung bei verspätetem Beitritt zu entrichten sind. Sie waren zum Teil sehr beträchtlich, wie die folgenden Beispiele zeigen. Doppelte Nachzahlungen haben geleistet:
\begin{center}
L. für 12 Jahre,\quad R. für 13 Jahre,\quad L. für 19 Jahre,\quad S. für 25 Jahre,\quad \mbox{C. für 30 Jahre.}
\end{center}
\noindent\hspace*{\parindent}Die Kommission ist daher der Meinung, dass ein so geringer und zweifelhafter Vorteil, wie er für die Kasse aus der Umwandlung des Instituts in eine Zwangsanstalt entstehen könnte, gegenüber der Zerstörung des bisherigen liberalen Charakters dieser Stiftung nicht ins Gewicht fallen darf.

Ein paar andere Mittel von ebenfalls nur schwacher oder unsicherer Wirkung werden am Schluss dieses Berichts erwähnt werden.

Als tatsächlich wirksame Mittel können demnach nur die folgenden betrachtet werden:

1) Erhöhung der Beiträge.

2) Herabsetzung der Pensionen.

3) Verbindung beider Mittel.

\Needspace{10\baselineskip}
Die in der Denkschrift aufgestellte Bilanzrechnung ermöglicht es, genau anzugeben, in welchem Maß diese Mittel angewandt werden müssen, wenn der Zweck erreicht, das heißt die Bilanz um 17520 Rthl. verbessert werden soll. Einzelheiten dazu wären hier nicht am Platz; das Endergebnis ist aber:

1) Soll das Defizit allein durch höhere Beiträge gedeckt werden, so müssen diese allgemein, das heißt für \textit{alle} jetzigen und künftigen Mitglieder, auf \(4\tfrac{1}{4}\) Louis d'or erhöht werden.

2) Soll die Herabsetzung der Pensionen allein die Deckung bewirken, so müssen diese auf 223 Rthl. verringert werden (wobei die Pension der Professorenwitwe G. bei 200 Rthl. bliebe).

3) Bei einer Verteilung der Last auf Beitragende und Pensionierte käme es darauf an, welches Teilungsverhältnis man wählte: Sollten zum Beispiel Erstere \(\tfrac{2}{3}\) und Letztere \(\tfrac{1}{3}\) übernehmen, so müssten die Beiträge \(3\tfrac{1}{2}\) Louis d'or und die Pensionen 241 Rthl. betragen. Sollen die Beiträge nur auf 3 Louis d'or erhöht werden, so können, wenn das Defizit wirklich gedeckt sein soll, nur 235 Rthl. Pension gezahlt werden.

Bei allen diesen Rechnungen ist vorausgesetzt, dass die gewählten Änderungen von Michaelis 1845 an wirksam werden und dass im zweiten oder dritten Fall die Pensionskürzungen sämtliche Witwen, die gegenwärtigen wie die künftigen, treffen. Sollten ohne jede Beitragserhöhung und ohne Pensionskürzung für die gegenwärtigen Witwen die künftigen Witwen allein die Last tragen, so könnte diesen nur eine Pension von \(213\tfrac{1}{4}\) Rthl. gewährt werden.

